\documentclass[aspectratio=169,10pt]{beamer}
\usetheme[numbering=fraction,progressbar=frametitle]{metropolis}
\usepackage{fontspec}
\setsansfont{Latin Modern Sans}
\usepackage{amsmath,amssymb,mathtools,graphicx,booktabs}
\usepackage{polyglossia}
\setdefaultlanguage{english}
\setbeamertemplate{navigation symbols}{}
\setbeamerfont{frametitle}{size=\large}
\setbeamerfont{normal text}{size=\normalsize}
\setlength{\parskip}{.45em}
\hypersetup{pdfauthor={},pdfsubject={Lecture slides},pdfkeywords={},pdftitle={Vectors in the plane}}
\providecommand{\N}{\mathbb{N}}
\providecommand{\Z}{\mathbb{Z}}
\providecommand{\D}{\mathbb{D}}
\providecommand{\Q}{\mathbb{Q}}
\providecommand{\R}{\mathbb{R}}
\providecommand{\C}{\mathbb{C}}
\providecommand{\vect}[1]{\overrightarrow{#1}}
\providecommand{\vi}{\vec{\imath}}
\providecommand{\vj}{\vec{\jmath}}
\providecommand{\vk}{\vec{k}}
\providecommand{\norm}[1]{\left\lVert #1\right\rVert}
\providecommand{\vc}[2]{\begin{pmatrix}#1\\#2\end{pmatrix}}
\providecommand{\deter}[2]{\det\left(#1,#2\right)}
\providecommand{\abs}[1]{\left\lvert #1\right\rvert}
\providecommand{\dd}{\,\mathrm{d}}
\providecommand{\e}{\mathrm{e}}
\providecommand{\Card}{\operatorname{Card}}
\providecommand{\E}{\operatorname{E}}
\providecommand{\V}{\operatorname{V}}
\providecommand{\degre}{^{\circ}}
\providecommand{\un}[1]{\,\mathrm{#1}}
\providecommand{\sci}[2]{#1\times10^{#2}}
\providecommand{\pgcd}{\operatorname{PGCD}}
\providecommand{\ppcm}{\operatorname{PPCM}}
\title{Vectors in the plane}
\subtitle{Lecture and worked examples}
\date{}
\author{}
\begin{document}
\begin{frame}[plain]
\vfill
{\usebeamerfont{title}\inserttitle\par}
\medskip
{\usebeamerfont{subtitle}\insertsubtitle\par}
\medskip
{\usebeamercolor[fg]{progress bar}\rule{\textwidth}{.5pt}}
\vfill
\end{frame}
\begin{frame}{Translation and vectors}
\begin{columns}[T,onlytextwidth]
\begin{column}{.57\textwidth}
The translation that transforms \(A\) into \(B\) associates with every point \(M\) the unique point \(M'\) such that \(ABM'M\) is a parallelogram, possibly flat. It is called the translation by the vector \(\vect{AB}\).
A nonzero vector has a direction, a sense and a norm. \(\|\vect{AB}\|=AB\). The vector \(\vect{AA}\) is the zero vector, written \(\vec0\); it has neither direction nor sense.
\end{column}
\begin{column}{.4\textwidth}
\centering
\includegraphics[width=\linewidth,height=4.5cm,keepaspectratio]{figures/seconde-02-1-en.pdf}
\end{column}
\end{columns}
\end{frame}
\begin{frame}{Translation and vectors}
\(\vect{AB}=\vect{CD}\) when the translation by the vector \(\vect{AB}\) transforms \(C\) into \(D\). Two nonzero vectors are equal if and only if they have the same direction, the same sense and the same norm.
The opposite vector of \(\vect{AB}\) is \(\vect{BA}\): \(\vect{BA}=-\vect{AB}\).
\end{frame}
\begin{frame}{Translation and vectors}
Draw a segment \([AB]\) and place a point \(C\) outside \((AB)\). Construct \(D\) such that \(\vect{AB}=\vect{CD}\). Name the parallelogram obtained. Construct \(E\) such that \(\vect{CE}=-\vect{AB}\).
\end{frame}
\begin{frame}{Worked answer: Translation and vectors}
\(D\) is constructed by carrying over from \(C\) the direction, sense and length of \(\vect{AB}\). The parallelogram is \(ABDC\), in that order. \(E\) is the reflection of \(D\) in the point \(C\), because \(\vect{CE}=-\vect{CD}\).
\end{frame}
\begin{frame}{Sum, difference and Chasles's relation}
\begin{columns}[T,onlytextwidth]
\begin{column}{.57\textwidth}
The sum \(\vec u+\vec v\) corresponds to carrying out the translations by the vectors \(\vec u\) and \(\vec v\) one after the other. For all points \(A,B,C\):
\[\boxed{\vect{AB}+\vect{BC}=\vect{AC}}\qquad\text{(Chasles's relation).}\]
We have \(\vec u-\vec v=\vec u+(-\vec v)\) and \(\vect{AB}+\vect{BA}=\vec0\).
\end{column}
\begin{column}{.4\textwidth}
\centering
\includegraphics[width=\linewidth,height=4.5cm,keepaspectratio]{figures/seconde-02-2-en.pdf}
\end{column}
\end{columns}
\end{frame}
\begin{frame}{Sum, difference and Chasles's relation}
\(ABDC\) is a parallelogram, possibly flat, if and only if \(\vect{AD}=\vect{AB}+\vect{AC}\).
\end{frame}
\begin{frame}{Sum, difference and Chasles's relation}
Simplify: \(\vect{AB}+\vect{BC}\); \(\vect{AB}+\vect{CA}\); \(\vect{AB}+\vect{BC}+\vect{CD}\); \(\vect{AB}-\vect{AC}\).
\end{frame}
\begin{frame}{Worked answer: Sum, difference and Chasles's relation}
\(\vect{AC}\); \(\vect{CA}+\vect{AB}=\vect{CB}\); \(\vect{AD}\); \(\vect{AB}+\vect{CA}=\vect{CB}\). The commutativity of addition allows us to reorder the vectors.
\end{frame}
\begin{frame}{Sum, difference and Chasles's relation}
In a parallelogram \(ABCD\), justify that \(\vect{AB}+\vect{AD}=\vect{AC}\).
\end{frame}
\begin{frame}{Worked answer: Sum, difference and Chasles's relation}
In \(ABCD\), \(\vect{AD}=\vect{BC}\). So \(\vect{AB}+\vect{AD}=\vect{AB}+\vect{BC}=\vect{AC}\).
\end{frame}
\begin{frame}{Product of a vector by a real number}
For \(k\in\R\), \(k\vec u=\vec0\) if \(k=0\) or \(\vec u=\vec0\). Otherwise, \(k\vec u\) has the same direction as \(\vec u\), the same sense if \(k>0\), the opposite sense if \(k<0\), and its norm is \(|k|\|\vec u\|\).
\[(k+\ell)\vec u=k\vec u+\ell\vec u,\quad k(\vec u+\vec v)=k\vec u+k\vec v,\quad k(\ell\vec u)=(k\ell)\vec u.\]
\end{frame}
\begin{frame}{Product of a vector by a real number}
For a vector \(\vec u\) of norm \(4\) cm, represent \(2\vec u\), \(-\vec u\) and \(\dfrac12\vec u\) from the same origin. State their norms.
\end{frame}
\begin{frame}{Worked answer: Product of a vector by a real number}
The norms are \(8\) cm, \(4\) cm, \(2\) cm. The vectors \(2\vec u\) and \(\dfrac12\vec u\) have the sense of \(\vec u\); \(-\vec u\) has the opposite sense.
\end{frame}
\begin{frame}{Product of a vector by a real number}
\(I\) is the midpoint of \([AB]\) if and only if \(\vect{AI}=\vect{IB}\), or equivalently \(\vect{AI}=\dfrac12\vect{AB}\), or \(\vect{IA}+\vect{IB}=\vec0\).
\end{frame}
\begin{frame}{Product of a vector by a real number}
In a triangle \(ABC\), \(I\) is the midpoint of \([AB]\) and \(J\) that of \([AC]\). Show that \(\vect{IJ}=\dfrac12\vect{BC}\).
\end{frame}
\begin{frame}{Worked answer: Product of a vector by a real number}
By Chasles, \(\vect{IJ}=\vect{IA}+\vect{AJ}=-\dfrac12\vect{AB}+\dfrac12\vect{AC}=\dfrac12(\vect{BA}+\vect{AC})=\dfrac12\vect{BC}\). We recover the parallelism of the lines \((IJ)\) and \((BC)\) when the triangle is not flat.
\end{frame}
\begin{frame}{Frames, bases and coordinates}
\begin{columns}[T,onlytextwidth]
\begin{column}{.57\textwidth}
Two non-collinear vectors \(\vec i,\vec j\) form a basis of the plane. Every vector \(\vec u\) can be written uniquely as \(\vec u=x\vec i+y\vec j\); its coordinates are \((x,y)\) or \(\vc{x}{y}\).
A frame \((O,\vec i,\vec j)\) associates with the point \(M\) the coordinates of \(\vect{OM}\). It is \textbf{orthonormal} if the basis vectors are perpendicular and of norm \(1\).
\end{column}
\begin{column}{.4\textwidth}
\centering
\includegraphics[width=\linewidth,height=4.5cm,keepaspectratio]{figures/seconde-02-3-en.pdf}
\end{column}
\end{columns}
\end{frame}
\begin{frame}{Frames, bases and coordinates}
If \(A(x_A,y_A)\), \(B(x_B,y_B)\), \(\vec u(x,y)\) and \(\vec v(x',y')\):
\[\vect{AB}=\vc{x_B-x_A}{y_B-y_A},\quad \vec u+\vec v=\vc{x+x'}{y+y'},\quad k\vec u=\vc{kx}{ky}.\]
Two vectors are equal if and only if their respective coordinates are equal.
\end{frame}
\begin{frame}{Frames, bases and coordinates}
For \(\vec u(3,2)\) and \(\vec v(-1,4)\), calculate \(\vec u+\vec v\) and \(2\vec u-\vec v\). Write \(\vec w(5,-3)\) in the basis \((\vec i,\vec j)\).
\end{frame}
\begin{frame}{Worked answer: Frames, bases and coordinates}
\(\vec u+\vec v=(2,6)\); \(2\vec u-\vec v=(7,0)\); \(\vec w=5\vec i-3\vec j\).
\end{frame}
\begin{frame}{Norm, distance and midpoint}
In an \textbf{orthonormal} frame, for \(\vec u(x,y)\) and two points \(A,B\):
\[\|\vec u\|=\sqrt{x^2+y^2},\qquad AB=\sqrt{(x_B-x_A)^2+(y_B-y_A)^2}.\]
In any frame, the midpoint of \([AB]\) is \(I\left(\dfrac{x_A+x_B}{2};\dfrac{y_A+y_B}{2}\right)\).
\end{frame}
\begin{frame}{Norm, distance and midpoint}
The distance formula follows from Pythagoras's theorem applied to the horizontal and vertical displacements. For the midpoint, \(\vect{AI}=\dfrac12\vect{AB}\) gives \(x_I-x_A=\dfrac{x_B-x_A}{2}\) and \(y_I-y_A=\dfrac{y_B-y_A}{2}\); we isolate \(x_I,y_I\).
\end{frame}
\begin{frame}{Norm, distance and midpoint}
In an orthonormal frame, \(A(3,0)\), \(B(5,3)\), \(C(4,0)\) and \(D(6,4)\).

\begin{enumerate}
\setcounter{enumi}{0}
\item Calculate \(\vect{AB}\), \(\vect{CD}\), then \(3\vect{AB}-5\vect{CD}\).

\item Determine \(U\) such that \(\vect{BU}=2\vect{AB}-\vect{CD}\).

\item Determine \(E\) such that \(ABCE\) is a parallelogram, then its centre \(I\).

\item Calculate \(AB\) and \(CD\).

\end{enumerate}
\end{frame}
\begin{frame}{Worked answer: Norm, distance and midpoint}
1. \(\vect{AB}=(2,3)\) and \(\vect{CD}=(2,4)\), so \(3\vect{AB}-5\vect{CD}=(-4,-11)\).

2. \(\vect{BU}=(2,2)\), so \(U=B+(2,2)=(7,5)\).

3. \(ABCE\) is a parallelogram if and only if \(\vect{EC}=\vect{AB}\), that is \((4-x_E,-y_E)=(2,3)\), so \(E(2,-3)\). The centre is the midpoint of \([AC]\): \(I(7/2,0)\). The midpoint of \([BE]\) gives the same result.

4. \(AB=\sqrt{13}\) and \(CD=\sqrt{20}=2\sqrt5\).
\end{frame}
\begin{frame}{Collinearity and determinant}
Two nonzero vectors are collinear if they have the same direction. The zero vector is collinear with every vector. If \(\vec v\ne\vec0\), \(\vec u\) and \(\vec v\) are collinear if and only if there exists a real number \(k\) such that \(\vec u=k\vec v\).
\end{frame}
\begin{frame}{Collinearity and determinant}
In the last characterisation, the condition \(\vec v\ne\vec0\) is necessary: a nonzero vector cannot be a multiple of the zero vector.
\end{frame}
\begin{frame}{Collinearity and determinant}
For \(\vec u(x,y)\), \(\vec v(x',y')\), \(\det(\vec u,\vec v)=xy'-yx'\).
\[\vec u,\vec v\text{ collinear}\quad\Longleftrightarrow\quad\det(\vec u,\vec v)=0.\]
\end{frame}
\begin{frame}{Collinearity and determinant}
If \(\vec u=k\vec v\), then \(xy'-yx'=kx'y'-ky'x'=0\). The case where a vector is zero is immediate.
Conversely, suppose \(\vec v\ne\vec0\) and \(xy'-yx'=0\). If \(x'\ne0\), let \(k=x/x'\). The equality gives \(y=ky'\), so \(\vec u=k\vec v\). If \(x'=0\), then \(y'\ne0\), hence \(x=0\); with \(k=y/y'\), we again obtain \(\vec u=k\vec v\).
\end{frame}
\begin{frame}{Collinearity and determinant}
Study the collinearity of: (a) \(\vec u(3,2)\) and \(\vec v(-6,-5)\); (b) \(\vec w(4,1)\) and \(\vec t(3,3/4)\). If \(\vect{AB}=2\vect{CD}\) and \(\vect{BD}=-5\vect{CD}\), relate \(\vect{AB}\) and \(\vect{BD}\).
\end{frame}
\begin{frame}{Worked answer: Collinearity and determinant}
(a) \(\det(\vec u,\vec v)=-15+12=-3\ne0\): not collinear. (b) \(\vec t=\dfrac34\vec w\): collinear; the determinant is \(4\times\dfrac34-3=0\). Finally \(\vect{AB}=-\dfrac25\vect{BD}\), including if the vectors are zero.
\end{frame}
\begin{frame}{Alignment and parallelism}
For \(A\ne B\) and \(C\ne D\), the lines \((AB)\) and \((CD)\) are parallel (possibly coincident) if and only if \(\det(\vect{AB},\vect{CD})=0\).
The points \(A,B,C\) are collinear if and only if \(\det(\vect{AB},\vect{AC})=0\).
\end{frame}
\begin{frame}{Alignment and parallelism}
Show that \((AB)\) and \((CD)\) are parallel with \(A(2,-4)\), \(B(10,-2)\), \(C(-16,-4)\) and \(D(12,3)\).
\end{frame}
\begin{frame}{Worked answer: Alignment and parallelism}
\(\vect{AB}=(8,2)\) and \(\vect{CD}=(28,7)\). Their determinant is \(8\times7-2\times28=0\). The points of each pair are distinct: both lines exist and are parallel.
\end{frame}
\begin{frame}{Alignment and parallelism}
Show that \(E(0,1/2)\), \(F(2,1)\) and \(G(6,2)\) are collinear.
\end{frame}
\begin{frame}{Worked answer: Alignment and parallelism}
\(\vect{EF}=(2,1/2)\) and \(\vect{EG}=(6,3/2)=3\vect{EF}\). The determinant is \(2\times3/2-(1/2)\times6=0\); so the points are collinear.
\end{frame}
\begin{frame}{Alignment and parallelism}
With \(A(-1,1)\), \(B(3,2)\), \(C(-2,-3)\), \(D(6,-1)\), \(E(5,0)\), show that \((AB)\parallel(CD)\) and that \(E,B,D\) are collinear.
\end{frame}
\begin{frame}{Worked answer: Alignment and parallelism}
\(\vect{AB}=(4,1)\) and \(\vect{CD}=(8,2)=2\vect{AB}\). Then \(\vect{EB}=(-2,2)\) and \(\vect{ED}=(1,-1)=-\dfrac12\vect{EB}\). Both results follow.
\end{frame}
\end{document}
